%% IEEE conference paper -- Overleaf-ready (compile with pdfLaTeX).
%% Working draft: verified results are final; full-stack integration is marked
%% in red as [PLACEHOLDER]/[TBD] and filled in as the integration lands.
%% Upload the 3 PNGs from experiments/results/ next to this file (optional:
%% it compiles without them via \IfFileExists).
\documentclass[conference]{IEEEtran}
\IEEEoverridecommandlockouts

\usepackage{cite}
\usepackage{amsmath,amssymb,amsfonts}
\usepackage{graphicx}
\usepackage{booktabs}
\usepackage{xcolor}
\usepackage{url}
\usepackage{algorithm}

% --- helpers ---
\newcommand{\phold}[1]{\textcolor{red}{[PLACEHOLDER: #1]}}
\newcommand{\tbd}{\textcolor{red}{[TBD]}}
\newcommand{\y}{\checkmark}
\newcommand{\n}{$\times$}
% figure-or-box so the paper compiles even before figures are uploaded
\newcommand{\figorbox}[2]{%
  \IfFileExists{#1}{\includegraphics[width=\columnwidth]{#1}}%
  {\fbox{\parbox[c][3.2cm][c]{0.92\columnwidth}{\centering \texttt{#2}}}}}
\newcommand{\figorboxwide}[2]{%
  \IfFileExists{#1}{\includegraphics[width=\textwidth]{#1}}%
  {\fbox{\parbox[c][2.6cm][c]{0.97\textwidth}{\centering \texttt{#2}}}}}

\begin{document}

\title{An Open, Site-Specific, Multi-Cell 5G Digital Twin:\\
Ray-Traced Throughput, an OpenAirInterface-Grounded\\
Link Model, and a Closed-Loop Traffic-Steering rApp}

\author{\IEEEauthorblockN{Author Name(s)}
\IEEEauthorblockA{\textit{Affiliation} \\
\textit{Organization}\\
City, Country \\
email@domain}}

\maketitle

\begin{abstract}
Network digital twins are increasingly used to develop and validate O-RAN control
applications before deployment, but open platforms trade off protocol-stack
realism, site-specific propagation, network scale, and closed-loop control---none
combines all four. We present an open-source, CPU-only digital twin that (i)~builds
a multi-cell, multi-user 5G scene from OpenStreetMap and ray-traces every
cell-to-user channel with NVIDIA Sionna~RT; (ii)~computes per-user downlink
throughput from a link model \emph{measured on OpenAirInterface's (OAI) real
physical layer}, using per-resource-element SINR with EESM link abstraction;
(iii)~hosts an O-RAN traffic-steering rApp; and (iv)~grounds per-link realism by
injecting the exact ray-traced channel into a live OAI gNB--UE link. Across three
cities and eight random layouts each (24 runs), the rApp raises median per-user
throughput by $9.5\pm11.3\%$ and served cell-edge throughput by $22.4\pm24.8\%$,
improves Jain fairness by $9.1\%$, and reduces peak cell load by $0.75$ users,
while leaving aggregate rate and coverage-limited outage unchanged. A configuration
sweep shows per-user throughput rising $1.5$--$1.7\times$ as inter-cell load falls
to idle. We further verify per-link real-stack grounding by injecting the exact
ray-traced channel into a live OAI link, and outline the path to a fully integrated
multi-UE closed loop. The platform runs on commodity CPU hardware and is released
open-source.
\end{abstract}

\begin{IEEEkeywords}
digital twin, 5G NR, OpenAirInterface, Sionna RT, ray tracing, O-RAN, rApp,
traffic steering, link abstraction.
\end{IEEEkeywords}

\section{Introduction}
The digital twin is emerging as the environment in which AI/ML-driven RAN control
(O-RAN rApps/xApps) is trained and validated before it touches a production
network, and its value hinges on data fidelity. Open tooling splits into camps that
do not meet: system-level simulators scale to many cells but abstract the physical
(PHY) layer; OAI-based emulators execute the real 5G-NR stack but, with realistic
propagation, cover a single link; commercial twins close the rApp loop but are
proprietary.

The fidelity that matters for control-application development is not merely a
faithful channel: it is the end-to-end mapping from site geometry to the KPIs a
controller observes and acts on. A traffic-steering rApp, for example, decides
which cell serves each user from signal, interference, and load; those quantities
depend jointly on 3D propagation (site-specific), on the neighbouring cells
(multi-cell), and on a link model that turns SINR into a real modulation-and-coding
outcome (protocol-stack realism). Evaluating such a controller therefore requires
all of these together, and---because random deployments vary widely---across many
layouts rather than a single scenario. This motivates both our platform and our
evaluation methodology.

Four capabilities matter \emph{simultaneously} for a useful RAN twin---a real
protocol stack, a site-specific ray-traced channel, a multi-cell network, and a
closed control loop---and no open platform combines them. This paper contributes
such a platform with a statistically rigorous, honest evaluation. Our
contributions are:
\begin{enumerate}
\item An open, CPU-only, multi-cell/multi-user twin whose throughput is grounded in
OAI's measured PHY, driven by per-resource-element SINR and EESM.
\item A traffic-steering rApp evaluated over 24 runs across three cities, with
robust statistics and explicit outage accounting.
\item Verified injection of the exact ray-traced channel into a live OAI link, plus
a twin-to-OAI inter-cell-interference noise-floor bridge, and a design for the
full-stack closed loop (multi-UE Standalone, FlexRIC KPM, and the rApp acting on
OAI) that these components enable.
\item A reproducible experiment harness and results, released open-source.
\end{enumerate}

\section{Background}
\textbf{Link-to-system abstraction.} System-level evaluation cannot afford full PHY
processing per user, so it relies on link-to-system (L2S) mapping: a link-level
simulator characterises the SINR-to-block-error-rate (BLER) relation for each MCS,
and an effective-SINR mapping collapses a frequency-selective SINR vector into a
single AWGN-equivalent value that indexes those curves. EESM is the widely used
one-parameter form. The accuracy of L2S hinges on how faithful the underlying
SINR-to-BLER curves are; we obtain ours from OpenAirInterface's real receiver
rather than from idealised expressions.

\textbf{O-RAN control loops.} The O-RAN architecture externalises RAN control into
applications: xApps on the near-real-time RIC (tens of ms, via E2) and rApps on the
non-real-time RIC (seconds and above, via O1/A1). Traffic steering and load
balancing are canonical rApp use cases, typically realised by adjusting per-cell
handover offsets (cell-individual offsets, CIO). A digital twin lets such an rApp be
trained and validated before it is allowed to touch a live network.

\textbf{Twin fidelity.} A twin used to rank or tune control policies must reproduce
not only average link quality but the joint distribution of the KPIs a controller
acts on, across the deployment randomness it will meet in the field. This makes
multi-scene, multi-layout evaluation---and honest reporting of variance and
outage---a fidelity requirement, not a nicety.

\section{Related Work}
Table~\ref{tab:related} positions the closest open efforts on the four axes.
Every existing open platform is missing at least one. OWDT~\cite{owdt} pairs OAI
with Sionna~RT for a faithful single gNB--UE link but is single-cell and
monitor-only. Tiny-Twin~\cite{tinytwin} runs a real OAI stack for multiple UEs with
a RIC loop, but single-cell from replayed channel traces (it lists inter-cell
interference as future work). Orange's RIC-TaaP/ns-O-RAN and OpenTwin~\cite{opentwin}
close rApp loops across many cells---OpenTwin adds twin-to-real KPM
self-correction---but on ns-3's abstracted PHY. Recent OAI+FlexRIC
work~\cite{oaiho} closes handover loops across real cells but over the air, without
a site-specific ray-traced channel. Within OAI, an emerging ray-tracing channel
emulator feeds the virtual-time simulator VRTSim (multi-node on its roadmap) and is
a natural backend our injection step can adopt. The asterisk on our fourth axis is
deliberate (Sec.~\ref{sec:disc}): the four are integrated as \emph{components} of
one platform, with the rApp loop evaluated on the analytical twin and the real
stack grounded per link; full rApp-driven multi-cell OAI is outlined as the next
step in Sec.~\ref{sec:results}-E.

Commercial network digital twins close this gap in proprietary form: Rimedo Labs'
NDT coupled with the VIAVI emulator trains and validates traffic-steering rApps
before deployment, reporting double-digit throughput gains. Such tools demonstrate
the industrial value of the closed loop but are neither open nor reproducible; our
aim is an open, commodity-hardware counterpart whose every layer---geometry,
propagation, link model, and control---can be inspected and re-run.

Our link model builds on the established link-to-system (L2S) abstraction used in
system-level simulation~\cite{eesm}: a link-level simulator characterises SINR-to-BLER
behaviour, and EESM compresses a frequency-selective SINR vector into an
AWGN-equivalent effective SINR for a fast network-scale lookup. The distinction
here is that our lookup is measured from OAI's own PHY (\texttt{nr\_dlsim}) rather
than from idealised curves, so the network-scale KPIs inherit a real receiver's
implementation losses. Our multi-cell interference-as-noise-floor bridge is, in
spirit, a site-specific generalisation of the 3GPP OFDMA channel noise generator
(OCNG)~\cite{ocng}, which loads unused resources to emulate other-cell activity;
we instead derive each UE's interference from the ray-traced neighbour channels.

\begin{table}[t]
\caption{Open RAN digital-twin platforms on four axes.}
\label{tab:related}
\centering
\footnotesize
\begin{tabular}{@{}lcccc@{}}
\toprule
Platform & Real & Site & Multi & Closed \\
 & stack & RT & cell & loop \\
\midrule
OWDT~\cite{owdt}                 & \y & \y & \n & \n \\
Tiny-Twin~\cite{tinytwin}        & \y & \n$^{a}$ & \n & \y \\
Orange ns-O-RAN / RIC-TaaP       & \n$^{b}$ & \y & \y & \y \\
OpenTwin~\cite{opentwin}         & \n$^{b}$ & \n & \y & \y \\
OAI+FlexRIC handover~\cite{oaiho}& \y & \n$^{c}$ & \y & \y \\
OAI VRTSim / RT emulator         & \y & \y & \n$^{d}$ & \n \\
\textbf{This work}               & \y & \y & \y & \y$^{*}$ \\
\bottomrule
\end{tabular}
\\[2pt]
\raggedright\scriptsize
$^{a}$replayed traces; $^{b}$ns-3 abstract PHY; $^{c}$over-the-air;
$^{d}$multi-node on roadmap; $^{*}$components integrated, full loop in progress.
\end{table}

\section{System Architecture}
The platform is a pipeline with two consumers of one physics substrate
(Fig.~\ref{fig:arch}): a fast analytical layer that covers every cell and UE, and a
real-stack layer that grounds sampled links in OAI. Both are driven by the same
ray-traced channel, and the network data source is swappable so that synthetic and
real deployments run through an identical path.

\begin{figure*}[t]
\centering
\figorboxwide{fig_architecture.png}{fig\_architecture.png : OSM $\to$ Sionna RT $\to$ \{analytical KPI engine + rApp\} and \{OAI real-stack grounding\}}
\caption{Architecture. OpenStreetMap geometry and a swappable network source feed
Sionna~RT; the ray-traced channel drives (top) an analytical multi-cell KPI engine
and the traffic-steering rApp, and (bottom) a live OAI link over the exact channel
with a per-UE interference noise floor.}
\label{fig:arch}
\end{figure*}

\textbf{Geometry.} A WGS84 bounding box drives an OpenStreetMap (Overpass) query;
footprints are extruded into a 3D Mitsuba scene with ITU materials.

\textbf{Network.} Cells and UEs come from a swappable source: a seeded random
generator ($N$ sites $\times$ 3 sectors, uniform UE drops) or a CSV source for real
coordinates (e.g., OpenCelliD).

\textbf{Channel.} Sionna~RT places one transmitter per cell (3GPP sector pattern,
planar array, downtilt) and one receiver per UE and computes the channel frequency
response for every cell-to-UE link over an OFDM resource grid.

\textbf{KPI engine.} Let $\mathbf{H}_{u,c}(k)\in\mathbb{C}^{N_r\times N_t}$ be the
ray-traced channel of link $c\!\to\!u$ on resource element (RE) $k$, and
$P^{\mathrm{tx}}_c$ the cell transmit power. The wideband received power is
$P^{\mathrm{rx}}_{u,c}=P^{\mathrm{tx}}_c\,\overline{\lVert\mathbf{H}_{u,c}\rVert_F^2}$,
where the bar denotes averaging over antennas and REs. Each UE associates to the
cell maximising received power biased by a cell-individual offset (CIO), the
O-RAN control knob:
\begin{equation}
s(u)=\arg\max_{c}\; P^{\mathrm{rx}}_{u,c}\cdot 10^{\,\mathrm{CIO}_c/10}.
\label{eq:assoc}
\end{equation}
Because each cell serves one UE at a time under airtime sharing, the optimal
single-stream precoder is the matched filter, so the serving per-RE gain is the
largest singular value $\sigma_{\max}$ of the serving channel. With flat per-RE
transmit power $p_c$, a neighbour load factor $\rho\in[0,1]$, and thermal noise
$N_0=kTB\cdot\mathrm{NF}$, the per-RE SINR is
\begin{equation}
\gamma_u(k)=\frac{p_{s}\,\sigma_{\max}\!\big(\mathbf{H}_{u,s}(k)\big)^2}
{\rho\sum_{c\neq s}\frac{p_c}{N_t}\lVert\mathbf{H}_{u,c}(k)\rVert_F^2+N_0}.
\label{eq:sinr}
\end{equation}
The per-RE SINRs are compressed to one effective SINR by exponential effective
SINR mapping (EESM), which preserves the frequency selectivity the ray tracer
produces:
\begin{equation}
\gamma^{\mathrm{eff}}_u=-\beta\,\ln\!\Big(\tfrac{1}{K}\textstyle\sum_{k}
e^{-\gamma_u(k)/\beta}\Big),
\label{eq:eesm}
\end{equation}
with per-modulation $\beta$ calibratable by a single scale factor. The effective
SINR is mapped to a modulation-and-coding scheme (MCS) and spectral efficiency
$\eta_u$ through a curve $\Phi_{\mathrm{OAI}}$ \emph{measured offline from OAI's
\texttt{nr\_dlsim}} (real LDPC decoding and rate matching), i.e.\
$(m_u,\eta_u)=\Phi_{\mathrm{OAI}}(\gamma^{\mathrm{eff}}_u)$. On a static
full-buffer snapshot, proportional-fair scheduling reduces to equal airtime, so
the per-UE throughput is
\begin{equation}
R_u=\frac{B_{s}}{K_{s}}\,\eta_u\,(1-\mathrm{BLER}_{\mathrm{tgt}}),
\label{eq:rate}
\end{equation}
where $B_s$ is the serving cell bandwidth and $K_s$ its number of attached UEs.

\textbf{Traffic-steering rApp.} The rApp adjusts the CIO vector to maximise the
proportional-fair utility $U(\mathrm{CIO})=\sum_u\log R_u$ (Alg.~\ref{alg:ts}).
Because $U$ is piecewise-constant in the CIOs (it changes only when a UE's serving
cell flips), a fixed small step stalls on plateaus; we therefore line-search each
cell's CIO over a grid (coordinate ascent), which is monotonically non-decreasing
in $U$ by construction. Raising a cell's CIO pulls boundary UEs in (absorbing
load); lowering it sheds them to neighbours.

\begin{algorithm}[t]
\caption{Proportional-fair traffic steering}
\label{alg:ts}
\footnotesize
\begin{enumerate}
\item $\mathrm{CIO}\leftarrow \mathbf{0}$; evaluate $R_u$ and $U$ via
Eqs.~\eqref{eq:assoc}--\eqref{eq:rate}.
\item \textbf{repeat} (sweep over cells):
  \begin{enumerate}
  \item for each cell $c$: line-search $\mathrm{CIO}_c$ over a grid (others fixed)
  and keep the value maximising $U$.
  \end{enumerate}
\item \textbf{until} no cell's line-search improves $U$ (local optimum).
\end{enumerate}
\end{algorithm}

\textbf{Real-stack grounding.} A bridge converts each link to time-domain taps
injected into a live OAI gNB--UE link (a patch to OAI's channel module); a second
bridge writes each UE's twin-computed inter-cell interference as a per-client OAI
noise floor, so a single-link OAI run reflects its multi-cell context.

\section{Experimental Setup}
\textbf{Scenes.} Berlin-Mitte (2{,}623 buildings), Paris-Op\'era (121),
Manhattan-Midtown (119); geometry fetched/built once per scene and reused across
seeds. \textbf{Random layouts.} Eight seeded layouts per scene (9 cells, 30 UEs)---
24 runs ($323$\,s on one CPU server). Radio: 3.5\,GHz, 20\,MHz,
46\,dBm, $4\times1$ BS array, NF 7\,dB, 30\,kHz SCS, ray-tracing
depth 3, default inter-cell load $1.0$. \textbf{Configuration sweep.} Neighbour
load factor in $\{0,0.25,0.5,0.75,1.0\}$ on a cached channel per scene.
\textbf{Metrics.} We report per-UE downlink throughput (median and mean); the
\emph{served cell-edge} rate, defined as the 5th percentile among UEs with non-zero
throughput (robust to coverage holes, unlike a raw 5th percentile which collapses
to zero once outage exceeds $5\%$); the \emph{outage rate}, the fraction of UEs with
zero throughput; Jain's fairness index over per-UE rates; and the peak cell load
(largest number of UEs on any cell). rApp gains are reported as mean $\pm$ standard
deviation over the eight seeds per scene, and pooled across all 24 runs; a gain is
undefined (excluded) for a run whose baseline metric is zero.

\section{Results}\label{sec:results}
\subsection{Traffic-steering rApp}
Across 24 runs, relative to strongest-cell association (Table~\ref{tab:results},
Fig.~\ref{fig:cdf}): the rApp helps the users a strongest-cell policy under-serves
(median, cell-edge, fairness) and sheds peak load. Three honest observations:
aggregate rate barely moves ($+0.5\%$; the proportional-fair objective trades sum
for fairness), outage is unchanged (steering relocates covered users, it does not
create coverage), and variance is high (layout-dependent)---motivating multi-seed
evaluation.

The per-scene breakdown is instructive. Berlin, the densest scene, shows the
largest served cell-edge gain ($+41\%$): dense blockage produces stronger
load imbalance for the rApp to exploit. Paris shows the smallest gain ($+6\%$)
and the highest outage ($23.8\%$): when many users lack coverage entirely,
steering has little covered load to rebalance. Manhattan sits in between. The
consistent signal across all three cities is that the direction of the effect is
robust---median throughput, cell-edge rate, and fairness improve while peak load
falls---whereas the \emph{magnitude} is strongly deployment-dependent, which is
exactly the kind of conclusion a single-scenario study would misrepresent.

\begin{table}[t]
\caption{rApp gains vs.\ strongest-cell (mean $\pm$ std over seeds).}
\label{tab:results}
\centering
\footnotesize
\begin{tabular}{@{}lcccc@{}}
\toprule
Metric & Overall & Berlin & Paris & Manh. \\
\midrule
Median tput gain    & $9.5\pm11.3\%$ & $10.5\%$ & $4.5\%$ & $13.3\%$ \\
Served edge gain    & $22.4\pm24.8\%$& $41.0\%$ & $6.0\%$ & $20.2\%$ \\
Jain fairness gain  & $9.1\pm11.0\%$ & $11.2\%$ & $4.4\%$ & $11.5\%$ \\
Peak load (UEs)     & $-0.75$        & $-0.88$  & $-0.50$ & $-0.88$ \\
Mean tput gain      & $0.5\pm5.8\%$  & $-0.7\%$ & $-0.3\%$& $2.5\%$ \\
Outage change       & $0.0$ pts      & $0.0$    & $0.0$   & $0.0$ \\
\bottomrule
\end{tabular}
\end{table}

\begin{figure}[t]
\centering
\figorbox{fig_cdf_pooled.png}{fig\_cdf\_pooled.png}
\caption{Per-UE throughput CDF, baseline vs.\ rApp (pooled over scenes/seeds).}
\label{fig:cdf}
\end{figure}

\subsection{Coverage and outage}
Mean UE outage is $16.2\%$ at full load, and scene-dependent (Berlin $8.8\%$,
Manhattan $16.3\%$, Paris $23.8\%$), reflecting random placement over the extent.
For a fixed layout, outage is coverage-limited (path-loss-driven) and insensitive
to inter-cell load, while served users are interference-limited
(Fig.~\ref{fig:outage}).

\begin{figure}[t]
\centering
\figorbox{fig_outage_vs_load.png}{fig\_outage\_vs\_load.png}
\caption{UE outage vs.\ inter-cell load factor, per scene.}
\label{fig:outage}
\end{figure}

\subsection{Sensitivity to inter-cell load}
As the neighbour load factor falls from $1.0$ to $0$, mean per-UE throughput rises
$1.5$--$1.7\times$ (Berlin $13.6\to22.7$, Paris $14.9\to25.2$, Manhattan
$11.9\to22.0$~Mbps; Fig.~\ref{fig:tput}), confirming served users are
interference-limited. Combined with the outage result, this separates two regimes
the twin must represent distinctly: a coverage-limited tail, set by path loss and
insensitive to load, that steering and load control cannot fix; and an
interference-limited body, where both the neighbour load and the rApp's
load-balancing act. A twin that reported only average throughput would blur these,
and an rApp tuned on such a twin could over-claim gains that evaporate in the
coverage-limited tail of a real deployment.

\begin{figure}[t]
\centering
\figorbox{fig_throughput_vs_load.png}{fig\_throughput\_vs\_load.png}
\caption{Mean per-UE throughput vs.\ inter-cell load factor, per scene.}
\label{fig:tput}
\end{figure}

\subsection{Real-stack grounding (single link)}
The bridge converts a link's ray-traced channel frequency response to a
sample-spaced impulse response (IFFT), retains the strongest causal taps,
normalises them to unit energy so their shape carries the multipath while the
overall gain is carried separately, and maps each tap delay to a sample index.
These taps override the software-radio-simulator's channel through a small patch to
OAI's channel module. A companion bridge expresses each UE's twin-computed
inter-cell interference as a noise-floor rise,
$\Delta N_u=10\log_{10}\!\big((I_u+N_0)/N_0\big)$, written as a per-client OAI noise
power, so a single-link run reflects its multi-cell context. OAI builds and runs on
the same CPU host; a gNB--UE link reports real PHY KPIs (downlink SNR
${\sim}15$\,dB, MCS 9), and injection is confirmed in the gNB log
(\texttt{[RT] injected 4 taps ... channel\_length=49}); the link then runs its
PHY/MAC (HARQ, link adaptation) over the site-specific channel. Absolute path loss
is masked by automatic gain control and uplink power control in phy-test mode, so
the proof of grounding is the injection and the multipath the receiver equalises,
not an absolute level change; absolute goodput requires the full-stack Standalone
path outlined in Sec.~\ref{sec:results}-E.

\subsection{Toward full-stack integration}
The verified components above compose into a full-stack closed loop, which we
describe as the immediate next step. A multi-UE Standalone deployment (OAI 5G core
in containers, gNB, and $N$ UEs in network namespaces) would place each UE on its
ray-traced channel and computed interference noise floor; FlexRIC would stream
per-UE KPMs (E2SM-KPM); iperf3 would supply real per-UE goodput; and the
traffic-steering rApp would issue CIO/handover actions over O1/telnet, closing the
loop on the real stack. The headline measurement this enables is the
\emph{fidelity gap}---the difference between the rApp gain the analytical twin
predicts (Table~\ref{tab:results}) and the gain the real stack delivers---arguably
the most useful single number a network digital twin can report. The bridges of
Sec.~III (channel injection and per-UE interference noise floor) are the mechanism;
the integration requires a host with Docker and root privileges.

\subsection{Implications for twin-based rApp development}
Two practical conclusions follow. First, a network digital twin intended to
\emph{rank} or \emph{tune} control policies should report distributions over many
random deployments, not a headline number from one scenario: our own single-seed
pilot suggested a uniformly large gain, which the 24-run study tempered into a
robust-direction/variable-magnitude result. Second, coverage and load are distinct
problems---traffic steering improves the experience of \emph{covered} users but
cannot rescue users in outage---so a twin used to size an rApp's benefit must
account for the coverage regime (here, via the outage rate and the load-sensitivity
sweep) rather than reporting served-user throughput alone. These are precisely the
kinds of design insights a twin is meant to surface before deployment.

\subsection{Reproducibility and open release}
All results are produced by an open experiment harness that fetches each scene's
geometry once, reuses it across seeds (so only the random placement and the ray
tracing re-run per layout), and aggregates per-run statistics with means, medians,
and outage. A single command reproduces the 24-run matrix and the load sweep on a
commodity CPU in about five minutes; each scene is ray-traced in roughly one
minute. The twin, the OAI-measured link curve $\Phi_{\mathrm{OAI}}$, the
traffic-steering rApp, and the OAI bridges (channel injection and per-UE
interference noise floor) are released together, so that both the fast-layer
network results and the per-link real-stack grounding can be re-run independently,
and the full-stack integration of Sec.~\ref{sec:results}-E can be completed on a
host with Docker and root.

\section{Discussion and Limitations}\label{sec:disc}
We state these plainly. (1)~The rApp closed loop is currently evaluated on the
multi-cell analytical twin; the OAI stack provides per-link ground truth and
verified channel injection but is not yet \emph{driven} by the rApp in one running
system. Section~\ref{sec:results}-E outlines this integration; the two bridges
(channel injection and interference noise floor) are the mechanism that connects
them. (2)~Inter-cell interference is modelled as (frequency-selective)
noise power rather than physically combined waveforms; waveform-level combining
across cells requires cluster-, FPGA-, or GPU-class channel emulators and is out of
scope for a CPU platform. (3)~The evaluation is a static full-buffer snapshot, all
cells share one carrier, and only the downlink is modelled; mobility (via a
virtual-time simulator), fractional frequency reuse, and uplink are natural
extensions. (4)~rApp benefits are high-variance and strongly layout-dependent, so
we report full distributions rather than a single figure; practitioners should
treat the mean as indicative and the spread as the operative uncertainty.
(5)~Traffic steering does not address coverage: users in outage cannot be rescued
by re-association, only by densification or power/tilt control. (6)~Validation is
stack-versus-model; calibration against field measurements, and the use of real
cell-site databases in place of random layouts, are the highest-value next steps.

\section{Conclusion}
We presented an open, CPU-only, site-specific, multi-cell 5G digital twin that
grounds throughput in OAI's physical layer, hosts a traffic-steering rApp, and
transmits a real OAI link over the exact ray-traced channel. Across three cities
and eight random layouts each, the rApp improves median and cell-edge throughput
and fairness and relieves peak load, without inflating aggregate rate or fixing
coverage---an honest, reproducible result. The full-stack closed loop
(Sec.~\ref{sec:results}-E) is the planned integration that will make it the first
open platform to bring all four ingredients together on commodity hardware.

\textbf{Future work.} Completing the full-stack loop (multi-UE Standalone, FlexRIC
KPM, and the rApp acting on OAI) will let us report the twin-predicted-vs-stack
fidelity gap on a control task---arguably the most useful single number a RAN twin
can provide. Beyond that, real cell-site databases in place of random layouts,
mobility through a virtual-time simulator, uplink and multi-carrier support, and
calibration against field measurements are the natural directions; each slots into
the existing pipeline without changing its structure.

\begin{thebibliography}{00}
\bibitem{owdt} T.~Iye \emph{et al.}, ``Open Wireless Digital Twin: End-to-End 5G
Mobility Emulation with OpenAirInterface and Ray Tracing,'' \emph{IEEE Access},
2025 (arXiv:2503.12177).
\bibitem{tinytwin} A.~Mamaghani \emph{et al.}, ``Tiny-Twin: A CPU-Native Full-stack
Digital Twin for NextG Cellular Networks,'' arXiv:2601.08217, 2026.
\bibitem{opentwin} ``OpenTwin: Digital Twin--Driven Closed-Loop KPM Inference and
Control for Open RAN,'' arXiv:2605.24662, 2026.
\bibitem{oaiho} ``Towards Intelligent Mobility Management: Customized Handover in 5G
O-RAN Networks,'' ACM, 2025.
\bibitem{sionnart} J.~Hoydis \emph{et al.}, ``Sionna RT: Differentiable Ray Tracing
for Radio Propagation Modeling,'' arXiv:2303.11103.
\bibitem{oai} F.~Kaltenberger \emph{et al.}, ``Driving Innovation in 6G Wireless
Technologies: The OpenAirInterface Approach,'' \emph{Computer Networks}, 2025.
\bibitem{flexric} R.~Schmidt, M.~Irazabal, and N.~Nikaein, ``FlexRIC: An SDK for
Next-Generation SD-RANs,'' in \emph{Proc. ACM CoNEXT}, 2021.
\bibitem{ns3sionna} R.~Pegurri \emph{et al.}, ``Toward Digital Network Twins:
Integrating Sionna RT in ns-3 for 6G Multi-RAT Networks,'' INFOCOM WKSHPS, 2025
(arXiv:2501.00372).
\bibitem{eesm} 3GPP TS 38.214, ``NR; Physical layer procedures for data.''
\bibitem{ocng} 3GPP TS 38.106, ``NR; Repeater radio transmission and reception''
(OCNG).
\end{thebibliography}

\end{document}
